\chapter{AI, Hukum, dan Manusia}
\label{ch:manusia}

Di sela kuliah teknis, saya mengikuti serangkaian mata kuliah yang melihat AI dari luar: AI and law,
yang diberikan oleh pengacara dan ahli hukum teknologi; AI in society dan technology and society dari
kacamata sosiologi; robopsychology dan communicating AI dari Martina Mara, yang laboratoriumnya meneliti
bagaimana manusia memandang dan memperlakukan mesin. Bab ini merangkum apa yang saya bawa pulang dari
kuliah-kuliah itu. Ia lebih sedikit rumusnya, tetapi pertanyaannya tidak kalah sulit.

\section{Apa itu sistem AI menurut hukum}

Hukum butuh definisi, dan definisi AI ternyata sulit. AI Act Uni Eropa \citep{aiact2024} memakai
definisi yang berakar pada pekerjaan kelompok ahli tingkat tinggi Komisi Eropa \citep{hleg2019}: sistem
berbasis mesin yang dirancang untuk beroperasi dengan berbagai tingkat otonomi, mungkin beradaptasi
setelah dipasang, dan untuk tujuan tertentu menyimpulkan dari input yang diterimanya cara menghasilkan
keluaran seperti prediksi, konten, rekomendasi, atau keputusan yang bisa memengaruhi lingkungan fisik
atau virtual. Kata kuncinya adalah \emph{menyimpulkan}. Program yang hanya menjalankan aturan yang ditulis
manusia secara eksplisit umumnya tidak termasuk, sedangkan model yang dilatih dari data termasuk.

\section{AI Act: regulasi berbasis risiko}

Slide kuliah menyusun linimasa AI Act: diusulkan Komisi Eropa pada April 2021, diterbitkan di Lembaran
Resmi Uni Eropa pada 12 Juli 2024, dan mulai berlaku pada 1 Agustus 2024, dengan ketentuan-ketentuannya
berlaku bertahap. Larangan untuk sistem dengan risiko yang tidak dapat diterima berlaku sejak Februari
2025, kewajiban bagi penyedia model AI serbaguna sejak Agustus 2025, dan menurut jadwal awal sebagian
besar ketentuan lainnya mulai Agustus 2026. Jadwal ini kemudian diubah. Paket penyederhanaan yang
dikenal sebagai Digital Omnibus on AI, yang disepakati pada Mei 2026 dan diterbitkan di Lembaran Resmi
pada akhir Juli 2026, menunda kewajiban untuk sistem berisiko tinggi yang berdiri sendiri ke 2 Desember
2027, dan untuk AI yang tertanam di produk yang sudah diatur regulasi lain ke 2 Agustus 2028. Kewajiban
transparansi, misalnya untuk konten sintetis, tetap berjalan sejak Agustus 2026, dengan masa transisi
sampai Desember 2026 untuk sistem yang sudah beredar. Ketika buku ini ditulis, peraturannya masih
bergerak, sehingga pembaca sebaiknya selalu memeriksa versi terbarunya.

Inti AI Act adalah piramida risiko. Di puncak ada praktik yang dilarang sama sekali: penilaian sosial
warga negara (\istilah{social scoring}), manipulasi yang mengeksploitasi kerentanan orang, pengenalan emosi
di tempat kerja dan sekolah, kategorisasi biometrik untuk menyimpulkan ras, keyakinan politik, atau
orientasi seksual, serta pengumpulan gambar wajah secara tidak tertarget. Di bawahnya ada sistem berisiko
tinggi, misalnya AI dalam rekrutmen, penilaian kredit, pendidikan, infrastruktur kritis, dan alat
kesehatan. Sistem ini boleh dipakai, tetapi penyedianya wajib menjalankan manajemen risiko, menjaga
kualitas data, membuat dokumentasi teknis, mencatat log, memastikan pengawasan manusia, serta memenuhi
syarat akurasi dan ketahanan. Lebih ke bawah, sistem dengan risiko terbatas seperti chatbot terutama
dikenai kewajiban transparansi: pengguna harus tahu bahwa mereka berbicara dengan mesin. Sistem lainnya
praktis bebas.

Model AI serbaguna, seperti model bahasa besar, ditambahkan menjelang akhir perundingan. Penyedianya
wajib menyiapkan dokumentasi teknis dan informasi bagi pihak yang membangun aplikasi di atasnya. Model
yang dilatih dengan total komputasi lebih dari $10^{25}$ operasi floating point dianggap berpotensi
membawa risiko sistemik dan dikenai kewajiban tambahan: evaluasi model, mitigasi risiko, pelaporan
insiden serius, dan keamanan siber. Untuk pelanggaran, denda bisa mencapai 35 juta euro atau tujuh persen
omzet global tahunan, mana yang lebih tinggi.

\section{Data pribadi dan GDPR}

Sebelum AI Act, perangkat hukum utama untuk AI di Eropa adalah GDPR \citep{gdpr2016}, yang berlaku sejak
25 Mei 2018. GDPR mengatur pemrosesan data pribadi, yaitu informasi apa pun yang berkaitan dengan orang
yang bisa diidentifikasi. Enam prinsipnya menjadi kerangka berpikir yang berguna bagi siapa pun yang
melatih model: pemrosesan harus sah, adil, dan transparan; data hanya dipakai untuk tujuan yang
dinyatakan; data yang dikumpulkan sesedikit mungkin; data harus akurat; data tidak disimpan lebih lama
dari perlu; dan integritas serta kerahasiaannya dijaga. Pengendali data bertanggung jawab membuktikan
kepatuhan pada semua prinsip itu.

Satu ketentuan berbicara langsung tentang AI. Seseorang berhak untuk tidak tunduk pada keputusan yang
semata-mata didasarkan pada pemrosesan otomatis, termasuk profiling, kalau keputusan itu berakibat
hukum atau berdampak signifikan baginya. Ada pengecualian, misalnya kalau diperlukan untuk kontrak,
diizinkan undang-undang, atau disetujui secara eksplisit, tetapi pesannya jelas: manusia harus tetap
punya tempat dalam keputusan penting. Bagi pengembang, ini berarti penilaian dampak perlindungan data
sebelum sistem berisiko dijalankan, dan perlindungan privasi yang dirancang sejak awal (\istilah{privacy by
design}) alih-alih ditambal belakangan.

Kuliah juga membahas pertanggungjawaban. Arahan tanggung jawab produk Uni Eropa yang direvisi memperluas
definisi produk sehingga mencakup perangkat lunak, termasuk sistem AI, dan meringankan beban pembuktian
bagi korban dalam kasus yang secara teknis rumit. Usulan arahan tanggung jawab AI membahas pengungkapan
bukti dan anggapan hubungan sebab-akibat yang bisa dibantah ketika sistem AI terlibat dalam kerugian.

\section{Teknologi dan masyarakat}

Kuliah sosiologi dimulai dengan konsep \istilah{sociological imagination} dari \citet{mills1959}: kemampuan
melihat bagaimana masalah pribadi terhubung dengan isu sosial yang lebih besar. Orang yang ditolak
kreditnya oleh algoritme mengalami masalah pribadi. Kalau ribuan orang dari lingkungan yang sama ditolak
dengan cara serupa, itu isu sosial.

Dari sana kuliah menolak \istilah{technological determinism}, pandangan bahwa teknologi berkembang menurut
logikanya sendiri lalu mengubah masyarakat. Sebaliknya, teknologi dibentuk oleh pilihan sosial: siapa yang
mendanai riset, data apa yang dikumpulkan, masalah siapa yang dianggap penting, dan bagaimana sebuah sistem
akhirnya dipakai. Strategi AI nasional berbagai negara menunjukkan bahwa satu teknologi yang sama
dibayangkan dengan berbeda di tempat yang berbeda. \citet{benjamin2019} memperlihatkan bagaimana
sistem yang tampak netral bisa mereproduksi ketidakadilan rasial, karena ia belajar dari data dan
keputusan masa lalu. \citet{zuboff2019} menyebut ekonomi yang mengubah perilaku manusia menjadi data untuk
memprediksi dan memengaruhinya sebagai kapitalisme pengawasan. Kita sudah melihat contoh kecilnya:
embedding kata yang menyerap stereotip gender \citep{bolukbasi2016}, sistem rekomendasi yang memperkuat
popularitas (\cref{ch:rekomendasi}), dan model medis yang belajar penanda rumah sakit alih-alih penyakit
(\cref{ch:medis}).

\section{Bagaimana manusia memandang mesin}

Robopsychology bukan psikologi untuk robot, melainkan psikologi manusia ketika berhadapan dengan mesin
yang tampak cerdas. Kuliahnya dimulai dengan metode: perbedaan antara psikologi sehari-hari dan psikologi
ilmiah, antara korelasi dan sebab-akibat, dan mengapa eksperimen dengan kelompok yang diacak dibutuhkan
untuk menyimpulkan efek. Laboratorium Mara, misalnya, memakai realitas virtual untuk menguji bagaimana
sinyal-sinyal sebuah robot kolaboratif memengaruhi perasaan aman pekerja di dekatnya.

Temuan klasik bidang ini adalah bahwa manusia memperlakukan komputer secara sosial hampir secara
otomatis. Orang bersikap sopan kepada komputer, menerapkan stereotip gender pada suara mesin, dan merasa
tersanjung oleh pujian dari program, meski tahu itu hanya program \citep{reeves1996,nass2000}.
\istilah{Anthropomorphism}, kecenderungan memberi sifat manusia pada benda, diperkuat oleh tiga faktor
menurut \citet{epley2007}: ketersediaan pengetahuan tentang manusia sebagai kerangka penjelasan,
kebutuhan untuk memahami dan meramalkan perilaku sesuatu, dan kebutuhan akan hubungan sosial. Semakin mirip
manusia sebuah robot, semakin ia disukai, sampai titik di mana kemiripan yang hampir tetapi tidak sempurna
justru terasa menyeramkan. Lembah ganjil ini, \istilah{uncanny valley}, diusulkan \citet{mori2012} pada 1970.
Kuliah juga membahas kepercayaan, kreativitas mesin, dan kebutuhan psikologis dasar manusia seperti
otonomi dan kompetensi ketika bekerja bersama teknologi.

\section{Berbicara tentang AI}

Kuliah communicating AI membahas literasi AI: kompetensi yang dibutuhkan orang awam untuk memahami,
menilai, dan memakai AI secara kritis \citep{long2020}. Kuliah memperlihatkan bahwa liputan media tentang
AI meningkat tajam sejak sekitar 2015, dan bahwa gambar yang dipakai media untuk mewakili AI didominasi
robot humanoid, otak bercahaya, dan warna biru. Sebuah survei yang dibahas di kuliah meminta peserta
menilai gambar-gambar seperti ini menurut seberapa mirip manusia, seberapa menyeramkan, dan seberapa
tepat mereka mewakili AI, lalu melihat bagaimana penilaian-penilaian itu saling berkaitan.
Narasi fiksi dan berita tentang AI cenderung berayun antara harapan berlebihan dan ketakutan akan
kepunahan, dan teori kultivasi menjelaskan bagaimana paparan berulang membentuk keyakinan publik.

Salah satu tugas kuliah meminta kami menganalisis artikel berita tentang AI, dan tugas kelompoknya
merancang kampanye literasi AI untuk kelompok sasaran tertentu. Pelajarannya relevan bagi siapa pun yang
bekerja dengan AI: cara kita menjelaskan sebuah sistem memengaruhi apakah orang terlalu percaya,
terlalu takut, atau memakainya dengan bijak.

\section{Apa artinya bagi pengembang}

Kuliah-kuliah ini meninggalkan beberapa kebiasaan yang patut dibawa ke pekerjaan teknis. Tanyakan sejak
awal siapa yang terdampak sistem yang dibangun dan bagaimana mereka bisa menggugat keputusannya. Periksa
kinerja model untuk setiap kelompok, bukan hanya rata-ratanya. Kumpulkan data seperlunya. Dokumentasikan
apa yang dilakukan model dan di mana batasnya, karena regulasi semakin mewajibkannya dan karena itu
adalah praktik teknik yang baik. Dan jelaskan sistem dengan jujur, tanpa melebih-lebihkan kecerdasannya
dan tanpa meremehkan risikonya.

\sumberbab{Bab ini mengikuti kuliah AI and Law I dan II (definisi AI, AI Act, GDPR, kontrak, tanggung
jawab produk), AI in Society dan Technology and Society (imajinasi sosiologis, pembentukan sosial
teknologi, AI dan ketidakadilan, kapitalisme pengawasan), Robopsychology (metode psikologi, antropomorfisme,
kepercayaan, kreativitas), dan Communicating AI (literasi AI, citra dan narasi AI di media).}
